\section{Bessel form of the scalar-return orbit sum}
\label{sec:scalar-return-bessel}
The phases below are the sixteen mode-pair phases. The canonical proof is
\path{mathematics/fixed-map/scalar_return_bessel.md}.
Integer comparisons are certified by
\path{formalization/Bgs/ScalarReturnBesselRational.lean}.
Python and C++ check the orbit counts and the Bessel sum in floating point.
No external refereeing is asserted.
BGS-TRACE-005 and BGS-BGS-001 remain open.


\begin{theorem}[Scalar return and the orbit Bessel sum, BGS-TRACE-129]
\label{thm:scalar-return-bessel}
Let \(N\ge 3\) be odd. If \(S^m\equiv I\pmod N\), then \(M_N^m\) is the scalar
\((\operatorname{Tr} M_N^m)/N\), and
\[
U_N(\eta,\varphi)^m
=\frac{\operatorname{Tr} M_N^m}{N}
\prod_{j=0}^{m-1}\exp(i\eta A_j(\varphi)),
\]
with \(A_j\) the Egorov transport of the phased kick. Modes congruent to \(0\)
modulo \(N\) contribute only a global phase. Their energy \(\sigma_0(N)\)
vanishes for every \(N>9\). For every odd \(N\), \(\det(S-I)=-4\), so
\(S-I\) is invertible and the order-two coefficient at time one vanishes:
the averaged squared trace equals \(1\) plus the remainder of BGS-TRACE-126.

For a position-diagonal kick, the classical orbit sum
\(\mathcal S_N\) has phase average
\[
\mathbb E\lvert\mathcal S_N\rvert^2
=N^{-4}\sum_{z,z'}\prod_r
J_0\bigl(2\eta\lvert c_r\rvert\,\lvert V_r(z)-V_r(z')\rvert\bigr).
\]
At an odd prime, if the first rows of \(S^0,\ldots,S^{m-1}\) are pairwise
distinct, each non-zero frequency has mean square \(m\), and the degree-two
jet of this average is \(1-m\sigma_{\mathrm{diag}}\eta^2\), the same jet as
\(\exp(-m\sigma_{\mathrm{diag}}\eta^2)\).

The exponential is not claimed as an identity for \(\eta^2\sigma m\) of
order one. The degree-four gap is the quadruple excess of the position
orbit. That excess vanishes at the primes \(7\), \(31\), \(41\) and \(47\),
and equals \(120\) at \(p=11\), where the Bessel average at
\(\lambda=1\) exceeds the exponential by a factor about \(1.23\). The
momentum row has energy \(9/160000\). No COE limit and no CUE crossover are
proved.
\end{theorem}
